% ============================================================
%  DOCUMENTO: Apuntes de Derecho Civil — Forma de los Actos
%             Jurídicos e Instrumentos
%  AUTOR: Isaac Edgar Camacho Ocampo
%  AÑO: 2024
%  Universidad de Buenos Aires (UBA)  |  Facultad de Derecho
% ============================================================

\documentclass[12pt,oneside]{book}

% ── METADATOS ───────────────────────────────────────────────
\newcommand{\universidad}{Universidad de Buenos Aires (UBA)}
\newcommand{\facultad}{Facultad de Derecho}
\newcommand{\carrera}{}
\newcommand{\materia}{Derecho Civil — Introducción}
\newcommand{\titulo}{Forma de los Actos Jurídicos e Instrumentos}
\newcommand{\autor}{Isaac Edgar Camacho Ocampo}
\newcommand{\anio}{2024}

% ── PAQUETES ────────────────────────────────────────────────
\usepackage[utf8]{inputenc}
\usepackage[T1]{fontenc}
\usepackage[spanish,es-nodecimaldot]{babel}

\usepackage[
    top=2.5cm,
    bottom=2.5cm,
    left=3cm,
    right=2.5cm,
    headheight=15pt
]{geometry}

\usepackage{amsmath}
\usepackage{amssymb}
\usepackage{mathtools}

\usepackage{graphicx}
\usepackage{float}
\usepackage{caption}
\usepackage{subcaption}

\usepackage{booktabs}
\usepackage{array}
\usepackage{longtable}
\usepackage{tabularx}

\usepackage{enumitem}

\usepackage{xcolor}
\usepackage[most]{tcolorbox}

\usepackage{fancyhdr}
\usepackage{ragged2e}

\usepackage[
    colorlinks=true,
    linkcolor=blue!50!black,
    citecolor=green!40!black,
    urlcolor=blue!60!black,
    pdftitle={\titulo},
    pdfauthor={\autor},
    pdfsubject={Apuntes universitarios},
    bookmarks=true
]{hyperref}

% ── RUTAS DE IMÁGENES ───────────────────────────────────────
\graphicspath{
    {./img/}
    {./graficos/}
    {./figuras/}
}

% ── ÍNDICE ──────────────────────────────────────────────────
\setcounter{tocdepth}{2}

% ── PÁRRAFOS ────────────────────────────────────────────────
\setlength{\parindent}{0pt}
\setlength{\parskip}{0.5em}

% ── LISTAS ──────────────────────────────────────────────────
\setlist[itemize]{topsep=0.3em, itemsep=0.2em}
\setlist[enumerate]{topsep=0.3em, itemsep=0.2em}

% ── VIUDAS Y HUÉRFANAS ──────────────────────────────────────
\clubpenalty=10000
\widowpenalty=10000

% ── ENCABEZADOS Y PIES ──────────────────────────────────────
\pagestyle{fancy}
\fancyhf{}
\fancyhead[L]{\nouppercase{\leftmark}}
\fancyhead[R]{\materia}
\fancyfoot[C]{\thepage}
\renewcommand{\headrulewidth}{0.4pt}
\renewcommand{\footrulewidth}{0pt}

\fancypagestyle{plain}{
    \fancyhf{}
    \fancyfoot[C]{\thepage}
    \renewcommand{\headrulewidth}{0pt}
}

% ── COMANDOS MATEMÁTICOS ────────────────────────────────────
\newcommand{\abs}[1]{\left|#1\right|}
\newcommand{\norm}[1]{\left\lVert#1\right\rVert}
\newcommand{\vect}[1]{\mathbf{#1}}
\newcommand{\deriv}[2]{\frac{d#1}{d#2}}
\newcommand{\pderiv}[2]{\frac{\partial#1}{\partial#2}}

% ── CAJAS ACADÉMICAS ────────────────────────────────────────
\newtcolorbox{definition}[1]{
    enhanced, breakable,
    title={Definición: #1},
    fonttitle=\bfseries,
    colback=blue!3,
    colframe=blue!50!black,
    boxrule=0.8pt, arc=2mm
}

\newtcolorbox{important}{
    enhanced, breakable,
    title={Importante},
    fonttitle=\bfseries,
    colback=yellow!8,
    colframe=orange!70!black,
    boxrule=0.8pt, arc=2mm
}

\newtcolorbox{example}[1]{
    enhanced, breakable,
    title={Ejemplo: #1},
    fonttitle=\bfseries,
    colback=green!4,
    colframe=green!40!black,
    boxrule=0.8pt, arc=2mm
}

\newtcolorbox{formula}[1]{
    enhanced, breakable,
    title={Disposición normativa: #1},
    fonttitle=\bfseries,
    colback=purple!4,
    colframe=purple!50!black,
    boxrule=0.8pt, arc=2mm
}

\newtcolorbox{exercise}[1]{
    enhanced, breakable,
    title={Ejercicio: #1},
    fonttitle=\bfseries,
    colback=gray!5,
    colframe=gray!60!black,
    boxrule=0.8pt, arc=2mm
}

\newtcolorbox{note}{
    enhanced, breakable,
    title={Nota},
    fonttitle=\bfseries,
    colback=white,
    colframe=gray!50,
    boxrule=0.6pt, arc=2mm
}

% ============================================================
\begin{document}
% ============================================================

% ── PORTADA ─────────────────────────────────────────────────
\begin{titlepage}
\thispagestyle{empty}
\begin{center}

\vspace*{1cm}

{\large \textsc{\universidad}}

\vspace{0.2cm}

{\large \textsc{\facultad}}

\vspace{3cm}

{\Huge \textbf{\materia}}

\vspace{1cm}

{\LARGE \titulo}

\vspace{0.8cm}

\textit{Comprender la forma en que el derecho organiza la voluntad es comprender el lenguaje del mundo jurídico.}

\vspace{1.2cm}

\rule{0.70\textwidth}{0.5pt}

\vspace{0.5cm}

{\large Apuntes de estudio}

\vspace{0.5cm}

\rule{0.70\textwidth}{0.5pt}

\vfill

{\large \autor}

\vspace{0.3cm}

{\large \anio}

\end{center}

\vfill

\begin{flushleft}
\textit{Este es un modesto aporte a los estudiantes de ``Derecho Civil --- Introducción'' no pretende sustituir el material oficial de la materia.}
\end{flushleft}

\end{titlepage}

% ── ÍNDICE ──────────────────────────────────────────────────
\tableofcontents
\newpage

% ============================================================
\chapter{La Forma del Acto Jurídico}
% ============================================================

\section{Concepto de forma}

\begin{definition}{Forma del acto jurídico}
La forma es la exteriorización de la voluntad de los sujetos respecto del objeto, en orden a conseguir el fin jurídico y su protección al celebrar el acto. Ningún acto puede carecer de forma; en este sentido, la forma es un elemento esencial de los actos jurídicos.
\end{definition}

La forma del acto jurídico puede ser muy variada. El ordenamiento contempla desde formas simples —como un gesto que expresa la voluntad de adquirir algo— hasta actos de gran solemnidad, como el matrimonio y el testamento, para los cuales la ley establece formas determinadas. Entre estos extremos, existen los contratos verbales, los contratos electrónicos celebrados mediante formularios por internet y los instrumentos privados y públicos.

La forma se vincula con dos cuestiones centrales del derecho civil:

\begin{itemize}
    \item La \textbf{validez} del acto: en ciertos casos, la forma es requisito de existencia del acto y su incumplimiento genera la nulidad.
    \item La \textbf{prueba} del acto: en otros casos, la forma es exigida por la ley con el único fin de facilitar la demostración de que el acto fue celebrado.
\end{itemize}

\section{Regulación en el Código Civil y Comercial}

El Código Civil y Comercial de la Nación (CCyCN) regula la forma del acto jurídico dentro del Capítulo Quinto dedicado a los actos jurídicos, en las siguientes secciones del Libro Primero, Título Cuarto:

\begin{itemize}
    \item \textbf{Sección Tercera:} Forma y prueba del acto jurídico.
    \item \textbf{Sección Cuarta:} Instrumento público.
    \item \textbf{Sección Quinta:} Escritura pública.
    \item \textbf{Sección Sexta:} Instrumentos privados y particulares.
\end{itemize}

\section{El formalismo en el derecho: ventajas y desventajas}

La regulación jurídica de las formas no es neutral. Exige ponderarse en función de sus ventajas y desventajas para el tráfico jurídico.

\subsection{Ventajas del formalismo}

\begin{table}[H]
\centering
\begin{tabularx}{\textwidth}{>{\bfseries\raggedright\arraybackslash}p{0.35\textwidth}
                              >{\raggedright\arraybackslash}X}
\toprule
Ventaja & Explicación \\
\midrule
Facilita la prueba del acto & El acto queda documentado y puede acreditarse con mayor certeza. \\[4pt]
Publicita el acto & El acto es conocido por los demás; ingresa al conocimiento público. \\[4pt]
Colabora en la percepción de impuestos & Como ocurre con los contratos de alquiler que deben informarse a la AFIP. \\[4pt]
Permite distinguir el acto concluido de los preparatorios & La forma cristaliza el acto jurídico. Los borradores de un testamento no son un testamento; éste queda concluido con la forma mandada por la ley. \\[4pt]
Protege a terceros & Los derechos adquiridos respecto de terceros resultan tutelados cuando se cumple la forma. El tercero conoce su situación jurídica; esto es especialmente relevante en materia de inmuebles. \\
\bottomrule
\end{tabularx}
\end{table}

\begin{example}{Testamento — cristalización del acto}
Una persona puede estar pensando en hacer un testamento, consultar a distintas personas y elaborar borradores. En un momento determinado, la forma cristaliza ese acto jurídico. Los borradores no sirven: no son un testamento. El testamento queda concluido con la forma mandada por la ley: por escritura pública o por instrumento ológrafo.
\end{example}

\subsection{Desventajas del formalismo}

Las formas más estrictas presentan también inconvenientes:

\begin{itemize}
    \item Las formas son más costosas y engorrosas.
    \item Existe mayor riesgo de que el acto resulte inválido por algún defecto formal.
    \item Le dan \textbf{menor celeridad} a las transacciones.
    \item Generan incomodidades indudables en el tráfico jurídico cotidiano.
\end{itemize}

Por estas razones, el ordenamiento jurídico impone formas más exigentes únicamente a los actos más importantes o trascendentes, donde la mayor seguridad justifica la mayor incomodidad.

% ============================================================
\chapter{Clasificación de los Actos Jurídicos según la Forma}
% ============================================================

\section{Forma esencial y forma instrumental}

Es necesario distinguir dos sentidos del concepto de forma:

\begin{itemize}
    \item \textbf{Forma esencial:} es la manifestación de la voluntad que realizan las partes para celebrar el acto jurídico.
    \item \textbf{Forma instrumental:} es cuando la ley impone que algunos actos deben celebrarse mediante documentos escritos, ya sean públicos o privados.
\end{itemize}

El Código Civil y Comercial se concentra en la forma instrumental. Los instrumentos podrán ser públicos o privados.

\section{Primera clasificación: actos formales y no formales}

\begin{definition}{Actos formales}
Son aquellos respecto a los cuales la ley impone el cumplimiento de ciertas formas o solemnidades que deben cumplirse para la validez del acto. Ejemplo: el testamento, el matrimonio.
\end{definition}

\begin{definition}{Actos no formales}
Son aquellos que no tienen ninguna solemnidad prescripta por la ley; las partes son libres de realizarlos con libertad de formas.
\end{definition}

\section{Segunda clasificación: actos solemnes y no solemnes}

Los actos formales, a su vez, se clasifican en solemnes y no solemnes:

\begin{definition}{Actos formales solemnes}
Son aquellos en los que el rigor formal impuesto por la ley se vincula con la validez misma del acto. Si no se cumple con la forma, el acto es inválido.
\end{definition}

\begin{definition}{Actos formales no solemnes (ad probationem)}
Son aquellos en los que la forma es prescripta por la ley únicamente para facilitar la prueba del acto, sin afectar su validez.
\end{definition}

\section{Solemnidad absoluta y solemnidad relativa}

Dentro de los actos formales solemnes, el Código establece una distinción fundamental:

\begin{table}[H]
\centering
\begin{tabularx}{\textwidth}{>{\raggedright\arraybackslash}X
                              >{\raggedright\arraybackslash}X}
\toprule
\textbf{Solemnidad absoluta} & \textbf{Solemnidad relativa} \\
\midrule
La solemnidad es parte esencial del acto. Si no se cumple, el acto es ineficaz por nulidad de forma y no produce ningún efecto. & El acto no es válido (no genera los efectos propios del acto), pero genera la obligación de otorgar la forma prescripta. \\[6pt]
\textit{Ejemplo:} El testamento ológrafo que no está escrito y firmado de puño y letra por el testador. & \textit{Ejemplo:} El boleto de compraventa de un inmueble. No produce el efecto de la transferencia de la propiedad, pero sí genera la obligación de escriturar. \\
\bottomrule
\end{tabularx}
\end{table}

\subsection{Conversión del acto jurídico}

\begin{formula}{Art. 285 CCyCN — Conversión del acto jurídico}
\textit{El acto que no se otorga en la forma exigida por la ley no queda concluido como tal, pero vale como acto en el que las partes se obligan a cumplir con la expresada formalidad, salvo que ella se exija bajo sanción de nulidad.}
\end{formula}

La \textbf{conversión del acto jurídico} opera en los supuestos de solemnidad relativa: el acto no produce los efectos propios del negocio jurídico que se procuraba celebrar, pero se convierte en otro acto —la obligación de otorgar el instrumento proyectado—. Cuando la forma es exigida bajo sanción de nulidad (solemnidad absoluta), el incumplimiento genera la nulidad y no produce ningún otro efecto.

\begin{example}{Boleto de compraventa de inmueble}
Una compraventa de inmueble celebrada mediante boleto de compraventa (instrumento privado) no produce el efecto traslativo de dominio, ya que la ley exige escritura pública para ello. Sin embargo, en virtud de la conversión del artículo 285, genera la obligación de las partes de escriturar, es decir, de otorgar el instrumento proyectado.
\end{example}

% ============================================================
\chapter{Principio de Libertad de Formas}
% ============================================================

\begin{formula}{Art. 284 CCyCN — Libertad de formas}
\textit{Si la ley no designa una forma determinada para la exteriorización de la voluntad, las partes pueden utilizar la que estimen conveniente. Las partes pueden convenir una forma más exigente que la impuesta por la ley.}
\end{formula}

El principio de libertad de formas es la regla general del sistema. Si la ley no designa una forma determinada, las partes son libres de elegir la que consideren conveniente. Esta libertad tiene dos dimensiones:

\begin{itemize}
    \item \textbf{Libertad de elección:} ante el silencio de la ley sobre la forma, las partes determinan libremente el modo de exteriorizar su voluntad.
    \item \textbf{Autonomía para elevar la forma:} las partes pueden convenir una forma más exigente que la impuesta por la ley. Por ejemplo, celebrar un contrato de locación por escritura pública, aunque la ley no lo requiera.
\end{itemize}

\begin{example}{Locación por escritura pública}
Las partes de un contrato de alquiler pueden convenir celebrarlo por escritura pública, aunque la ley no exija esa forma para la locación. Ello implica costos adicionales (honorarios del escribano), pero resulta perfectamente válido.
\end{example}

Por supuesto, cuando la ley manda una forma determinada, las partes deben partir de la forma legalmente impuesta. El artículo 284 es el punto de partida para aquellos actos respecto de los cuales la ley no designa forma especial.

% ============================================================
\chapter{Los Instrumentos — Clasificación General}
% ============================================================

\section{Concepto de instrumento}

El concepto de instrumento refiere a la forma escrita. Es la expresión de la voluntad a través de un documento que da mayor seguridad y certeza acerca de los derechos y obligaciones asumidos por las partes.

\begin{formula}{Art. 286 CCyCN — Expresión escrita}
\textit{La expresión escrita puede tener lugar por instrumentos públicos, o por instrumentos particulares firmados o no firmados. Puede hacerse constar en cualquier soporte, siempre que su contenido sea representado con texto inteligible, aunque su lectura exija medios técnicos.}
\end{formula}

\begin{important}
La referencia a medios técnicos contempla los soportes informáticos y electrónicos. El Código extiende la categoría de instrumento a todo soporte inteligible, independientemente del medio requerido para su lectura.
\end{important}

\section{Clasificación de los instrumentos}

\begin{table}[H]
\centering
\begin{tabularx}{\textwidth}{>{\raggedright\arraybackslash}p{0.30\textwidth}
                              >{\raggedright\arraybackslash}p{0.30\textwidth}
                              >{\raggedright\arraybackslash}X}
\toprule
\textbf{Instrumentos públicos} & \textbf{Instrumentos privados} & \textbf{Particulares no firmados} \\
\midrule
Interviene un oficial público. Tienen validez y fuerza probatoria por sí mismos con los alcances que les concede el Código. & Surgen de las partes. Son instrumentos particulares firmados. La \textbf{firma} es un requisito esencial. & Todo escrito no firmado: impresos, registros visuales o auditivos de cosas o hechos, registros de la palabra y de la información, cualquiera sea el medio empleado. \\
\bottomrule
\end{tabularx}
\end{table}

\begin{formula}{Art. 287 CCyCN — Instrumentos particulares}
\textit{Los instrumentos particulares pueden estar firmados o no. Si lo están, se llaman instrumentos privados. Si no lo están, se los denomina instrumentos particulares no firmados; esta categoría comprende todo escrito no firmado, entre otros, los impresos, los registros visuales o auditivos de cosas o hechos y, cualquiera que sea el medio empleado, los registros de la palabra y de la información.}
\end{formula}

El Código Civil y Comercial amplió el tratamiento jurídico de esta materia respecto del Código anterior, incorporando la categoría de instrumentos particulares no firmados para dar cobertura a la realidad del tráfico contemporáneo, caracterizado por la informatización: formularios web, constancias electrónicas, mensajes de texto, registros digitales.

\begin{example}{Constancia de compra por internet}
Una constancia generada por una plataforma de comercio electrónico al completar un formulario de compra en línea —que registra los datos del contrato, el precio y el método de pago— es un instrumento particular no firmado. No tiene firma en sentido estricto, por lo que no es instrumento privado con su fuerza probatoria plena, pero puede tener valor como prueba en caso de conflicto.
\end{example}

\section{La firma como disposición general}

\begin{formula}{Art. 288 CCyCN — Firma}
\textit{La firma prueba la autoría de la declaración de voluntad expresada en el texto al cual corresponde. La firma consiste en el nombre del firmante o en un signo; en ambos casos se realiza con el trazo habitual y de tal manera que identifique a la persona.}
\end{formula}

En el Código Civil y Comercial, la regulación de la firma fue trasladada de la sección de instrumentos privados —donde se encontraba en el Código de Vélez— a las disposiciones generales sobre la forma, lo que significa que rige tanto para los instrumentos públicos (incluyendo la escritura pública) como para los privados. La firma:

\begin{itemize}
    \item Prueba la \textbf{autoría} de la declaración de voluntad.
    \item Expresa la \textbf{conformidad} con el contenido del documento.
    \item Consiste en el nombre o en un signo habitual e identificatorio de la persona.
\end{itemize}

\subsection{Firma digital y firma electrónica}

El artículo 288 contempla también la cuestión de los instrumentos generados por medios electrónicos y remite a la normativa específica sobre firma digital.

\begin{important}
La \textbf{firma digital} no debe confundirse con el agregado de datos personales al pie de un correo electrónico (nombre, cargo, teléfono). Este último no constituye firma digital en sentido jurídico.
\end{important}

\begin{definition}{Firma digital (Ley 25.506 — año 2001)}
Se entiende por firma digital al resultado de aplicar a un documento digital un procedimiento matemático (algoritmo) que requiere información de exclusivo conocimiento del firmante y que se encuentra bajo su absoluto control. La firma digital debe haber sido: (a) creada durante el período de vigencia del certificado digital válido del firmante; (b) verificada por referencia a los datos de verificación de firma; y (c) generada por un certificado emitido por un certificador licenciado.
\end{definition}

\begin{definition}{Firma electrónica}
Es un conjunto de datos electrónicos, ligados o asociados de manera lógica a otros datos electrónicos, que el signatario utiliza con fines de identificación, pero que no alcanza el estándar de la firma digital. No cuenta con certificación indubitable de autoría e integridad.
\end{definition}

La diferencia entre ambas figuras es fundamental en materia probatoria:

\begin{table}[H]
\centering
\begin{tabularx}{\textwidth}{>{\raggedright\arraybackslash}X
                              >{\raggedright\arraybackslash}X}
\toprule
\textbf{Firma digital} & \textbf{Firma electrónica} \\
\midrule
Certificación digital indubitable de \textbf{autoría} e \textbf{integridad} mediante tercero certificador licenciado. El documento firmado digitalmente no puede ser modificado con posterioridad sin que ello resulte detectable. & Identificación electrónica sin certificación indubitable. Puede ser desconocida; quien la invoca debe acreditar su validez mediante prueba. \\[6pt]
\textbf{Mayor valor probatorio.} Garantizadas la autoría y la integridad. & \textbf{Menor valor probatorio.} La autoría y la integridad no están garantizadas por estas certificaciones. \\
\bottomrule
\end{tabularx}
\end{table}

% ============================================================
\chapter{Instrumentos Públicos}
% ============================================================

\section{Concepto y caracteres}

\begin{definition}{Instrumento público}
Son los documentos a los cuales la ley les asigna un valor de autenticidad por sí mismos y que, como tales, no requieren de un reconocimiento previo de la firma para hacer plena fe de su contenido. La autenticidad del instrumento público proviene de la ley, no de la voluntad de las partes.
\end{definition}

En la mayoría de los instrumentos públicos interviene un oficial público —funcionario público o escribano— que certifica la validez del acto. Sin embargo, lo que otorga la fuerza probatoria propia al instrumento público es la ley: es ella la que, ante la intervención de determinado funcionario revestido de esa credibilidad social e institucional, concede al instrumento el carácter de auténtico.

Los caracteres esenciales del instrumento público son:

\begin{itemize}
    \item \textbf{Autenticidad:} hace plena fe de su contenido sin necesidad de reconocimiento previo.
    \item \textbf{Fuerza probatoria propia:} su valor probatorio proviene de la ley.
    \item \textbf{Oponibilidad erga omnes:} tiene valor frente a todos, no sólo entre las partes.
\end{itemize}

\section{Enumeración legal}

\begin{formula}{Art. 289 CCyCN — Enumeración de instrumentos públicos}
\textit{Son instrumentos públicos: a) Las escrituras públicas y sus copias o testimonios; b) Los instrumentos que extienden los escribanos o los funcionarios públicos con los requisitos que establecen las leyes; c) Los títulos emitidos por el Estado nacional, provincial o la Ciudad Autónoma de Buenos Aires, conforme a las leyes que autorizan su emisión.}
\end{formula}

Cada categoría comprende:

\begin{enumerate}
    \item \textbf{Escrituras públicas y sus copias o testimonios:} instrumento matriz del protocolo notarial y sus reproducciones autenticadas.

    \item \textbf{Instrumentos de escribanos o funcionarios públicos:} comprende, entre otros, las actas notariales, las sentencias judiciales homologadas, las actas celebradas en sede judicial, y los convenios celebrados ante funcionario competente (como un convenio de alimentos labrado en sede judicial).

    \item \textbf{Títulos emitidos por el Estado:} títulos de deuda pública nacional, provincial o de la Ciudad Autónoma de Buenos Aires, emitidos conforme a las leyes que autorizan su emisión.
\end{enumerate}

\section{Requisitos de validez}

\begin{formula}{Art. 290 CCyCN — Requisitos de validez del instrumento público}
\textit{Son requisitos de validez del instrumento público: a) La actuación del oficial público en los límites de sus atribuciones y de su competencia territorial, excepto que el lugar sea generalmente tenido como comprendido en ella; b) Las firmas del oficial público, de las partes, y en su caso, de sus representantes; si alguno de ellos no firma por sí o no puede hacerlo, otro de los intervinientes lo hará a su ruego, haciéndose constar esta circunstancia.}
\end{formula}

Los dos requisitos esenciales son:

\begin{itemize}
    \item \textbf{Competencia del oficial público:} comprende la competencia \emph{funcional} (atribuciones para celebrar ese tipo de acto) y la competencia \emph{territorial} (actuación dentro de la jurisdicción asignada). La actuación fuera de estos límites invalida el instrumento.
    \item \textbf{Firmas:} deben firmar el oficial público, las partes y, en su caso, los representantes. La falta de firma de cualquiera de los intervinientes invalida el instrumento para todos.
\end{itemize}

\begin{important}
Es presupuesto de validez que el oficial público se encuentre efectivamente en funciones. Los actos otorgados antes de la notificación de la suspensión o cesación de funciones son válidos. Si la persona ejercía efectivamente un cargo existente bajo apariencia de legitimidad, la falta de investidura no afecta al acto ni al instrumento dentro de los límites de la buena fe.
\end{important}

\section{Defectos de forma}

Los defectos de forma en el instrumento público —enmiendas, agregados, borraduras, entrelineados, alteraciones— deben ser salvados al final del instrumento antes de las firmas. Si no se salvan, el instrumento carece de validez en esa parte.

\begin{example}{Enmienda en número de DNI}
Si en una escritura pública se consignó erróneamente el número del documento nacional de identidad de uno de los comparecientes y se corrigió mediante una enmienda manuscrita sobre el texto impreso, esa corrección debe ser salvada al final del instrumento antes de las firmas. De lo contrario, la enmienda no salvada afecta la validez del instrumento.
\end{example}

\begin{note}
Si el instrumento contiene las firmas de las partes pero adolece de un defecto que le impide valer como instrumento público (por ejemplo, por incompetencia del escribano), puede igualmente valer como \textbf{instrumento privado}, en la medida en que el acto de que se trate no exija forma pública bajo sanción de nulidad.
\end{note}

\section{Eficacia probatoria del instrumento público}

Este es uno de los aspectos de mayor relevancia práctica. El instrumento público hace plena fe, pero el Código distingue dos situaciones según el tipo de contenido:

\begin{formula}{Art. 296 CCyCN — Eficacia probatoria del instrumento público}
\textit{El instrumento público hace plena fe: a) En cuanto a que se ha realizado el acto, la fecha, el lugar y los hechos que el oficial público enuncia como cumplidos por él o ante él hasta que sea declarado falso en juicio civil o criminal; b) En cuanto al contenido de las declaraciones sobre convenciones, disposiciones, pagos, reconocimientos y enunciaciones de hechos directamente relacionados con el objeto principal del acto instrumentado, hasta que se produzca prueba en contrario.}
\end{formula}

\subsection{Hechos percibidos directamente por el oficial público}

Los \textbf{hechos que el oficial público enuncia como cumplidos por él o ante él} —la fecha, el lugar, la celebración del acto, los hechos ocurridos en su presencia— hacen plena fe y sólo pueden ser cuestionados mediante un \textbf{juicio de redargución de falsedad} (acción civil o criminal tendiente a demostrar la falsedad del instrumento). Quien pretenda impugnarlos debe acusar al oficial público de haber faltado a la verdad, lo que supone una acusación de suma gravedad dada la credibilidad que la sociedad deposita en la función notarial.

\begin{example}{Acta notarial — constancia de hechos ante el escribano}
Si un escribano labra un acta en la que consigna: ``En el día de la fecha me constituí en el domicilio de la calle X y constaté que se realizaba la actividad Y, en presencia de las personas Z'', quien pretenda desacreditar ese hecho debe iniciar un juicio de redargución de falsedad, dado que estaría afirmando que el escribano mintió en su función oficial.
\end{example}

\subsection{Declaraciones de las partes}

En cuanto al \textbf{contenido de las declaraciones de las partes} —convenciones, disposiciones, pagos, reconocimientos y enunciaciones de hechos directamente relacionados con el objeto del acto—, el instrumento hace plena fe pero \textbf{admite prueba en contrario}. En este caso, no es necesario iniciar un juicio de redargución de falsedad: basta con demostrar por otros medios que el contenido de esas declaraciones no es veraz.

\begin{example}{Reconocimiento de deuda simulado}
Si en una escritura pública una de las partes reconoce adeudar una suma de dinero, ese reconocimiento es el contenido de su declaración de voluntad. Si posteriormente se pretende demostrar que ese reconocimiento fue simulado, no es necesaria la redargución de falsedad —el escribano no mintió: certificó que el compareciente hizo esa declaración—, sino que puede acreditarse la simulación por prueba en contrario.
\end{example}

% ============================================================
\chapter{La Escritura Pública}
% ============================================================

\section{Concepto y protocolo notarial}

\begin{formula}{Art. 299 CCyCN — Escritura pública. Definición}
\textit{La escritura pública es el instrumento matriz extendido en el protocolo de un escribano público o de otro funcionario autorizado para ejercer las mismas funciones, que contiene uno o más actos jurídicos.}
\end{formula}

La escritura pública es el tipo de instrumento público más importante en la regulación civil y comercial. Sus características esenciales son:

\begin{itemize}
    \item \textbf{Instrumento matriz:} la escritura original es aquella que queda integrada al protocolo del escribano. Lo que las partes conservan en sus domicilios es una \textbf{copia} o \textbf{testimonio}, que también tiene carácter de instrumento público y hace plena fe tanto como la escritura matriz.
    \item \textbf{Contenido:} la escritura contiene \textbf{actos jurídicos}. Esto la distingue del acta notarial, cuyo objeto es la comprobación de \textbf{hechos}.
    \item \textbf{Función notarial:} los escribanos son regulados en cada provincia y en la Ciudad Autónoma de Buenos Aires. El Código Civil y Comercial establece reglas comunes para este instrumento con validez en todo el país.
\end{itemize}

\begin{important}
Si existe discordancia entre el texto de la copia o testimonio y el de la escritura matriz, prevalece el contenido de la escritura matriz.
\end{important}

\subsection{El protocolo notarial}

El protocolo está conformado por folios numerados correlativamente. No pueden quedar hojas en blanco —lo que implicaría la posibilidad de antedatar documentos. Los asientos deben realizarse en orden cronológico. La reglamentación del libro del protocolo es de competencia local (provincial o de la Ciudad Autónoma de Buenos Aires).

\section{Requisitos de la escritura pública}

La escritura pública debe ser recibida personalmente por el escribano, quien tiene el \textbf{deber de calificación} de todos los presupuestos y elementos del acto: verifica la identidad de los comparecientes, su capacidad, la legitimidad del acto y la regularidad de los títulos.

\subsection{Estructura: encabezado, cuerpo y pie}

El contenido de la escritura pública (artículo 305 CCyCN) se organiza en tres grandes partes:

\begin{table}[H]
\centering
\begin{tabularx}{\textwidth}{>{\bfseries\raggedright\arraybackslash}p{0.20\textwidth}
                              >{\raggedright\arraybackslash}X}
\toprule
Parte & Contenido \\
\midrule
Encabezado & Lugar y fecha de otorgamiento (con posibilidad de incluir la hora). Datos de los comparecientes: nombre, apellido, documento nacional de identidad, domicilio real, estado civil. En el caso de segundas nupcias, se hace constar. Nombre del cónyuge, relevante por cuestiones vinculadas a la vivienda y al régimen patrimonial. Si comparece una persona jurídica: denominación completa, domicilio y datos del representante y de su mandato. \\[4pt]
Cuerpo & Naturaleza del acto e individualización de los bienes que constituyen su objeto. En las transmisiones de derechos reales sobre inmuebles, incluye el \textbf{``corresponde''}: la relación de antecedentes de titularidad del bien transmitido. \\[4pt]
Pie / Cierre & Constancia de la lectura del instrumento por el escribano ante los comparecientes. Salvado de enmiendas, borrones o entrelineados realizados antes de las firmas. Firma de los otorgantes, del escribano y de los testigos, si los hubiera. Si algún compareciente no puede firmar, otro firmará a su ruego, dejándose constancia del impedimento y de la impresión digital del otorgante. \\
\bottomrule
\end{tabularx}
\end{table}

\subsection{El ``corresponde'' en las transmisiones de inmuebles}

El \textbf{``corresponde''} es una cláusula obligatoria en las escrituras que instrumentan transmisiones de derechos reales sobre inmuebles. Contiene la relación de todos los antecedentes en virtud de los cuales el transmitente llegó a ser titular del derecho que transmite.

\begin{example}{Cadena de antecedentes}
Si el transmitente adquirió el inmueble por herencia, el corresponde consigna: ``Le corresponde por declaratoria de herederos de fecha X y partición de fecha Y. A su vez, el causante lo había adquirido por escritura de fecha Z, que obra en el folio N del Registro de Propiedad Inmueble.'' Esta cláusula permite encadenar las distintas escrituras y verificar la titularidad del bien a lo largo del tiempo.
\end{example}

La omisión del ``corresponde'' no genera la invalidez de la escritura, pero puede acarrear sanciones para el escribano en el ejercicio de su función profesional.

\section{Copias y testimonios}

Las copias o testimonios de la escritura matriz deben entregarse a las partes. Pueden reproducirse por cualquier medio que asegure su permanencia. Si alguna de las partes requiere una nueva copia, debe entregársele, salvo que exista una obligación de dar dinero pendiente de cumplimiento.

% ============================================================
\chapter{Actas Notariales}
% ============================================================

\begin{formula}{Art. 310 CCyCN — Actas notariales}
\textit{Se denominan actas los documentos notariales que tienen por objeto la comprobación de hechos.}
\end{formula}

\begin{definition}{Acta notarial}
Es el documento expedido por un escribano a requerimiento de parte, cuyo objeto es la comprobación de \textbf{hechos}. A diferencia de la escritura pública —que contiene actos jurídicos (Art. 299)—, el acta sólo constata hechos, sin crear, modificar ni extinguir relaciones jurídicas.
\end{definition}

El acta notarial implica la concurrencia del escribano a un lugar determinado para constatar un hecho y dejar constancia de él. Entre sus aplicaciones habituales se encuentran:

\begin{itemize}
    \item Constatación del estado de un bien al momento de su restitución (locaciones).
    \item Intimaciones de pago al deudor.
    \item Comprobación de que un acreedor se negó a recibir un pago.
    \item Constancia de hechos relevantes para el ejercicio de derechos.
\end{itemize}

\begin{example}{Acta de constatación al recuperar un local comercial}
Al recuperar la posesión de un local comercial después de una locación, el locador puede requerir al escribano que concurra al acto de entrega de llaves y labre un acta notarial que constate el estado en que se encuentra el inmueble. El escribano describe en el acta todo lo que constata, y ese documento tendrá valor probatorio en caso de reclamarse daños o el incumplimiento de las obligaciones de restitución.
\end{example}

Los requisitos del acta notarial están regulados en el artículo 311 CCyCN. El acta tiene valor probatorio respecto de:

\begin{itemize}
    \item Los \textbf{hechos} que el notario tuvo a la vista: su existencia, características y estado.
    \item Las \textbf{personas} que concurrieron, si quisieron identificarse.
    \item Las \textbf{declaraciones} efectuadas por los presentes.
\end{itemize}

\begin{note}
La Ciudad Autónoma de Buenos Aires regula el ejercicio de la función notarial mediante la Ley 404 (año 2000). El escribano puede ser titular o adscripto. El ingreso al notariado requiere la aprobación de un procedimiento de selección competitivo y una investidura formal.
\end{note}

% ============================================================
\chapter{Instrumentos Particulares: Privados y No Firmados}
% ============================================================

\section{Clasificación y jerarquía}

El Código Civil y Comercial sistematiza los instrumentos particulares en dos categorías:

\begin{itemize}
    \item \textbf{Instrumentos particulares firmados} (instrumentos privados): requieren firma de las partes como requisito esencial.
    \item \textbf{Instrumentos particulares no firmados}: todo escrito no firmado, registro visual o auditivo, o cualquier otro soporte inteligible sin firma.
\end{itemize}

La categoría más importante es la de los instrumentos privados (particulares firmados). Los instrumentos particulares no firmados constituyen una categoría incorporada por el nuevo Código para dar cobertura jurídica a la realidad del tráfico contemporáneo de bienes y servicios.

\section{Instrumentos privados: el requisito de la firma}

\begin{important}
El \textbf{requisito esencial} de los instrumentos privados es la \textbf{firma} de las partes (Art. 288). Sin firma, el documento es un instrumento particular no firmado, con menor fuerza probatoria.
\end{important}

El Código Civil y Comercial eliminó el requisito del \textbf{doble ejemplar} que imponía el Código anterior (Código de Vélez). En virtud de ello, ya no es jurídicamente exigible que los instrumentos se celebren en tantos ejemplares como partes intervinientes. No obstante, la buena práctica profesional recomienda mantener este recaudo, especialmente cuando una de las partes asume obligaciones relevantes.

\section{Reconocimiento de la firma}

El instrumento privado tiene una firma que prueba la declaración de voluntad, pero para que se produzca la plenitud de la fuerza probatoria entre las partes, la firma debe ser \textbf{reconocida}.

\begin{formula}{Art. 314 CCyCN — Reconocimiento de la firma}
\textit{Todo aquel contra quien se presente un instrumento cuya firma se le atribuye debe manifestarse al respecto. El silencio se tendrá por reconocimiento de la firma.}
\end{formula}

\begin{important}
El silencio frente a la intimación de reconocimiento de firma equivale al reconocimiento. Esto constituye un supuesto de silencio como manifestación de la voluntad (Art. 263 CCyCN): ante la obligación legal de expedirse, el silencio importa aceptación tácita.
\end{important}

El reconocimiento de la firma puede ser:

\begin{itemize}
    \item \textbf{Expreso:} la persona declara que la firma le pertenece.
    \item \textbf{Tácito:} guarda silencio ante la intimación, o realiza gestiones sin desconocer la autoría de la firma.
    \item \textbf{Forzoso:} se acredita mediante prueba, especialmente a través de la \textbf{pericia caligráfica} (cotejo de firmas realizado por un perito calígrafo designado judicialmente).
\end{itemize}

\begin{example}{Negación de contrato de locación}
Si el locatario niega haber suscripto el contrato de alquiler, el locador puede exhibirle el instrumento en juicio e intimarlo a que reconozca la firma. Si la niega, podrá acreditarse la autoría mediante pericia caligráfica, que establecerá con alta probabilidad la autenticidad de la firma atribuida al demandado.
\end{example}

El reconocimiento de la firma importa el reconocimiento de todo el cuerpo del instrumento: las cláusulas, disposiciones, convenciones y reconocimientos que contiene. Lo que el instrumento público tiene por sí, el instrumento privado lo adquiere con el reconocimiento de firma.

\section{Firma a ruego e impresión digital}

\begin{formula}{Art. 313 CCyCN — Firma a ruego e impresión digital}
\textit{Si alguno de los firmantes no sabe o no puede firmar, puede dejarse constancia de la impresión digital o bien de la firma de otra persona a ruego, en presencia de dos testigos que también lo suscriban.}
\end{formula}

Cuando una parte no sabe o no puede firmar, el Código prevé dos alternativas:

\begin{enumerate}
    \item \textbf{Firma a ruego:} otra persona firma en nombre del que no puede hacerlo, en presencia de dos testigos que también suscriben el instrumento. Esta modalidad tiene pleno valor probatorio, similar a un mandato.

    \item \textbf{Impresión digital:} la estampación de la huella dactilar del signatario. Sin embargo, conforme al artículo 314 CCyCN, el documento signado con impresión digital vale únicamente como \textbf{principio de prueba por escrito} y puede ser impugnado. Su menor valor probatorio obedece a que la impresión digital puede obtenerse sin el pleno consentimiento de la persona.
\end{enumerate}

\section{Documento firmado en blanco}

\begin{formula}{Art. 315 CCyCN — Documento firmado en blanco}
\textit{El firmante de un documento en blanco puede impugnar su contenido mediante la prueba de que no responde a sus instrucciones, pero no puede valerse de testigos si no existe principio de prueba por escrito. El desconocimiento del documento no puede invocarse para afectar al tercero que de buena fe celebró un acto jurídico basándose en él.}
\end{formula}

El documento firmado en blanco implica una suerte de mandato implícito: quien firma en blanco autoriza al receptor del documento a completarlo dentro de las instrucciones acordadas. Si el documento fue completado en forma contraria a esas instrucciones, el firmante puede impugnar su contenido, pero con importantes limitaciones probatorias:

\begin{itemize}
    \item No puede valerse exclusivamente de testigos si no existe principio de prueba por escrito.
    \item Si el instrumento fue utilizado de buena fe por un tercero que celebró un acto jurídico basándose en él, ese tercero queda protegido: el firmante no puede invocar el desconocimiento del documento en su perjuicio.
\end{itemize}

\begin{important}
La firma de documentos en blanco es una práctica de alto riesgo, ya que quien tiene el documento puede completarlo con obligaciones enormemente perjudiciales para el patrimonio del firmante.
\end{important}

\section{Enmiendas en instrumentos privados}

Las enmiendas, raspaduras o entrelineados que afecten partes esenciales del instrumento privado deben ser salvadas al final, con la firma de las partes. Diferencia con la escritura pública: en los instrumentos privados, las enmiendas no salvadas no generan la nulidad del instrumento, sino que el juez determinará en qué medida el defecto excluye o reduce su fuerza probatoria.

% ============================================================
\chapter{Fecha Cierta y Oponibilidad a Terceros}
% ============================================================

\section{El instrumento privado frente a terceros}

El instrumento privado tiene fuerza probatoria entre las partes desde el reconocimiento de la firma. Frente a terceros, sin embargo, rige una regla diferente: el instrumento privado \textbf{no es oponible} a terceros mientras no adquiera \textbf{fecha cierta}.

\begin{formula}{Art. 317 CCyCN — Fecha cierta}
\textit{La eficacia probatoria de los instrumentos privados reconocidos se extiende a los terceros desde su fecha cierta. Adquieren fecha cierta el día en que acontece un hecho del que resulta como consecuencia ineludible que el documento ya estaba firmado o no pudo ser firmado después.}
\end{formula}

\begin{definition}{Fecha cierta}
No es la fecha consignada en el instrumento, sino el momento a partir del cual existe certeza objetiva de que ese instrumento ya estaba firmado o no pudo ser firmado con posterioridad. Sirve para hacer oponible el instrumento privado a terceros.
\end{definition}

\section{Medios para adquirir fecha cierta}

El Código de Vélez (artículo 1035) enumeraba taxativamente los casos de adquisición de fecha cierta:

\begin{itemize}
    \item Exhibición del instrumento en juicio o ante una repartición pública donde quedara archivado.
    \item Reconocimiento ante un escribano y testigos.
    \item Transcripción en cualquier registro público.
    \item Fallecimiento de la parte que lo firmó.
\end{itemize}

El Código Civil y Comercial \textbf{abandona la enumeración taxativa} y establece que la fecha cierta puede acreditarse por cualquier medio de prueba, que el juez apreciará de manera rigurosa.

\begin{example}{Fecha cierta por exhibición en juicio}
Cuando un instrumento privado se incorpora a un expediente judicial, la mesa de entradas de tribunales le estampa un sello que acredita la fecha de presentación. Desde ese momento existe certeza de que el instrumento ya existía: puede oponerse a terceros como instrumento firmado con anterioridad a esa fecha.
\end{example}

\begin{example}{Fecha cierta por fallecimiento del firmante}
Si uno de los firmantes de un instrumento privado fallece, desde ese momento se tiene certeza de que el instrumento fue firmado antes de la muerte: con posterioridad al fallecimiento, esa persona ya no pudo haberlo suscripto. El instrumento adquiere fecha cierta desde el día del fallecimiento.
\end{example}

% ============================================================
\chapter{Instrumentos Particulares no Firmados y Correspondencia}
% ============================================================

\section{Concepto y amplitud}

Los instrumentos particulares no firmados comprenden todos los escritos no firmados, impresos, registros visuales o auditivos de cosas o hechos, cualquiera sea el medio empleado, y los registros de la palabra y de la información. Esta amplísima categoría abarca tanto los soportes en papel como los medios electrónicos: mensajes de texto, correos electrónicos, publicaciones en redes sociales, tickets electrónicos, registros de plataformas web.

\begin{note}
A diferencia de los instrumentos privados, los instrumentos particulares no firmados no cuentan con firma. La firma prueba la autoría de la declaración de voluntad; su ausencia implica que el instrumento tiene menor fuerza probatoria y que la valoración de su alcance queda librada a la apreciación judicial.
\end{note}

\section{Fuerza probatoria — criterios del juez}

\begin{formula}{Art. 319 CCyCN — Valor probatorio del instrumento particular}
\textit{El valor probatorio de los instrumentos particulares debe ser apreciado por el juez ponderando, entre otras pautas, la congruencia entre lo sucedido y narrado, la precisión y claridad técnica del texto, los usos y prácticas del tráfico, las relaciones precedentes y la confiabilidad de los soportes utilizados y de los procedimientos técnicos que se apliquen.}
\end{formula}

Las pautas que el juez debe ponderar son:

\begin{enumerate}
    \item \textbf{Congruencia entre lo sucedido y lo narrado:} correspondencia entre lo que dice el instrumento y los hechos ocurridos y acreditados por otros medios.

    \item \textbf{Precisión y claridad técnica del texto:} un formulario web completo y preciso tendrá mayor valor que un mensaje informal con redacción ambigua.

    \item \textbf{Usos y prácticas del tráfico:} la forma habitual en que se concluyen los actos jurídicos en ese sector de actividad o negocio.

    \item \textbf{Relaciones precedentes entre las partes:} el historial de intercambios previos que contextualice el instrumento en cuestión.

    \item \textbf{Confiabilidad de los soportes y procedimientos técnicos empleados:} la seguridad e integridad de los sistemas informáticos utilizados para generar o transmitir el instrumento.
\end{enumerate}

\begin{example}{Ticket electrónico como prueba}
En un conflicto sobre el pago del precio en una compraventa de bienes, el comprador presenta el ticket generado por la plataforma de comercio electrónico y el resumen de tarjeta de crédito que registra el débito correspondiente. La congruencia entre ambos elementos es un indicio de peso que el juez puede ponderar para tener por acreditado el pago.
\end{example}

\section{La correspondencia}

\begin{formula}{Art. 318 CCyCN — Correspondencia}
\textit{La correspondencia, cualquiera sea el medio empleado para crearla o transmitirla, puede presentarse como prueba por el destinatario, pero la que es confidencial no puede ser utilizada sin consentimiento del remitente.}
\end{formula}

La correspondencia —cartas, correos electrónicos, mensajes— pertenece, una vez recibida, al destinatario. Éste puede utilizarla como prueba en juicio contra el remitente, con la excepción de la \textbf{correspondencia confidencial}: si la comunicación tiene carácter confidencial, no puede ser utilizada sin el consentimiento del remitente.

La confidencialidad puede surgir:
\begin{itemize}
    \item De la manifestación expresa de las partes (por ejemplo, que el remitente la haya designado así).
    \item De la naturaleza misma de la comunicación (una carta dirigida a un amigo de confianza sobre asuntos personales).
\end{itemize}

Esta protección tiene fundamento constitucional: la \textbf{inviolabilidad de la correspondencia y los papeles privados} reconocida por la Constitución Nacional como garantía vinculada a la intimidad y privacidad de las personas.

\begin{example}{Uso de correspondencia confidencial}
Si alguien envía una carta a un amigo en la que reconoce adeudar sumas de dinero a terceros, y ese amigo entrega la carta a los acreedores para que la usen en un juicio, el uso de esa correspondencia puede ser cuestionado por su carácter confidencial. Sin el consentimiento del remitente, los destinatarios originales no pueden utilizarla.
\end{example}

% ============================================================
% CUADRO CORNELL
% ============================================================

\chapter{Cuadro Cornell — Repaso Integral}

\begin{table}[H]
\centering
\begin{tabularx}{\textwidth}{
    >{\raggedright\arraybackslash}p{0.30\textwidth}
    >{\raggedright\arraybackslash}X
}
\toprule
\textbf{Pregunta / Concepto clave} &
\textbf{Notas / Respuesta} \\
\midrule

\textbf{¿Qué es la forma del acto jurídico? ¿Es un elemento esencial?} &
Es la exteriorización de la voluntad de los sujetos respecto del objeto, en orden a conseguir el fin jurídico. Ningún acto puede carecer de forma: en ese sentido, la forma es un elemento esencial. \\[6pt]

\textbf{¿Cuáles son las dos dimensiones de la forma?} &
La forma tiene incidencia sobre la \textbf{validez} del acto (en ciertos casos, su incumplimiento produce nulidad) y sobre la \textbf{prueba} del acto (en otros casos, la forma es exigida sólo para facilitar su demostración). \\[6pt]

\textbf{¿Cuáles son las ventajas y desventajas del formalismo?} &
\textit{Ventajas:} facilita la prueba, publicita el acto, colabora en la percepción impositiva, permite distinguir el acto concluido de los preparatorios y protege a los terceros. \textit{Desventajas:} mayor costo, mayor riesgo de invalidez por defecto formal, menor celeridad en las transacciones, incomodidades prácticas. \\[6pt]

\textbf{¿Cuál es el principio general que rige la forma? (Art. 284)} &
El \textbf{principio de libertad de formas}: si la ley no designa una forma determinada, las partes pueden utilizar la que estimen conveniente. Las propias partes pueden pactar formas más exigentes que las legales. \\[6pt]

\textbf{¿Cómo se clasifican los actos según la forma?} &
Actos \textbf{formales} (con forma legal impuesta) y \textbf{no formales} (libertad de formas). Los formales se subdividen en \textbf{solemnes} (la forma es requisito de validez) y \textbf{no solemnes} o \emph{ad probationem} (la forma es exigida sólo para la prueba). \\[6pt]

\textbf{¿Qué diferencia hay entre solemnidad absoluta y solemnidad relativa?} &
En la \textbf{solemnidad absoluta}, el incumplimiento produce la nulidad del acto, sin ningún otro efecto. En la \textbf{solemnidad relativa}, el acto no produce los efectos propios del negocio, pero genera la obligación de otorgar la forma prescripta (conversión del acto jurídico, Art. 285). \\[6pt]

\textbf{¿Qué es la conversión del acto jurídico? (Art. 285)} &
Es el efecto propio de la solemnidad relativa: el acto no concluido en la forma requerida por la ley no queda concluido como tal, pero vale como acto por el cual las partes se obligan a cumplir con la formalidad, siempre que ésta no sea exigida bajo sanción de nulidad. Ejemplo: el boleto de compraventa genera la obligación de escriturar. \\[6pt]

\textbf{¿Cuáles son las categorías de instrumentos en el CCyCN? (Arts. 286--287)} &
(1) Instrumentos \textbf{públicos}: con autenticidad propia por ley. (2) Instrumentos \textbf{privados}: particulares firmados; requisito esencial: firma (Art. 288). (3) Instrumentos \textbf{particulares no firmados}: todo escrito, registro visual, auditivo o electrónico sin firma. \\[6pt]

\textbf{¿Qué prueba la firma? ¿Qué diferencia hay entre firma digital y firma electrónica? (Art. 288, Ley 25.506)} &
La firma prueba la \textbf{autoría} de la declaración de voluntad. La \textbf{firma digital} garantiza autoría e integridad mediante algoritmo matemático y certificado emitido por certificador licenciado. La \textbf{firma electrónica} es sólo un conjunto de datos identificatorios sin esa certificación; puede ser desconocida y su validez debe acreditarse. \\[6pt]

\textbf{¿Qué son los instrumentos públicos y cuáles son sus caracteres? (Art. 289)} &
Son los documentos a los que la ley les atribuye autenticidad propia, sin necesidad de reconocimiento previo. Hacen plena fe de su contenido. Enumeración: escrituras públicas y copias, instrumentos de escribanos o funcionarios públicos, y títulos emitidos por el Estado. \\[6pt]

\textbf{¿Cuáles son los requisitos de validez del instrumento público? (Art. 290)} &
(a) Actuación del oficial público dentro de sus \textbf{atribuciones} (competencia funcional) y \textbf{competencia territorial}. (b) \textbf{Firmas} del oficial, las partes y sus representantes. La falta de cualquier firma invalida el instrumento para todos. \\[6pt]

\textbf{¿Cómo funciona la eficacia probatoria del instrumento público? (Art. 296)} &
Distingue dos regímenes: (a) Los \textbf{hechos percibidos por el oficial} (fecha, lugar, actos cumplidos ante él) sólo pueden cuestionarse mediante \textbf{juicio de redargución de falsedad}. (b) El \textbf{contenido de las declaraciones de las partes} hace plena fe pero admite \textbf{prueba en contrario} sin necesidad de redargución. \\[6pt]

\textbf{¿Qué es la escritura pública y en qué se diferencia del acta notarial? (Arts. 299 y 310)} &
La \textbf{escritura pública} es el instrumento matriz extendido en el protocolo notarial que contiene \textbf{actos jurídicos}. El \textbf{acta notarial} es el documento notarial que tiene por objeto la comprobación de \textbf{hechos}. La distinción es esencial: la escritura crea, modifica o extingue derechos; el acta sólo constata hechos. \\[6pt]

\textbf{¿Qué estructura tiene la escritura pública? (Art. 305)} &
Tres partes: (1) \textbf{Encabezado}: lugar, fecha y datos de los comparecientes. (2) \textbf{Cuerpo}: naturaleza del acto, individualización de bienes y, en transmisiones de inmuebles, el ``corresponde''. (3) \textbf{Pie}: constancia de lectura, salvado de enmiendas, firmas. \\[6pt]

\textbf{¿Qué es el ``corresponde'' y para qué sirve?} &
Es la cláusula que, en las escrituras de transmisión de derechos reales sobre inmuebles, relaciona todos los antecedentes de titularidad del bien. Permite encadenar las sucesivas escrituras y verificar la historia dominial. Su omisión no invalida la escritura pero puede generar sanciones para el escribano. \\[6pt]

\textbf{¿Cuál es el requisito esencial del instrumento privado y cómo se reconoce la firma? (Arts. 288 y 314)} &
El requisito esencial es la \textbf{firma}. Para que el instrumento tenga plena fuerza probatoria entre las partes, la firma debe ser \textbf{reconocida}: expresa, tácita (el silencio frente a la intimación equivale a reconocimiento, Art. 314) o forzosa (pericia caligráfica). El reconocimiento de firma implica el reconocimiento de todo el contenido del instrumento. \\[6pt]

\textbf{¿Qué es la firma a ruego y cuál es la diferencia con la impresión digital? (Art. 313)} &
La \textbf{firma a ruego} es la firma de un tercero en nombre de quien no sabe o no puede firmar, en presencia de dos testigos. Equivale a la firma y otorga plena fuerza probatoria. La \textbf{impresión digital} sólo vale como \emph{principio de prueba por escrito} (Art. 314) y puede ser impugnada, por su menor seguridad. \\[6pt]

\textbf{¿Qué es la fecha cierta y para qué sirve? (Art. 317)} &
Es el momento a partir del cual existe certeza objetiva de que el instrumento ya estaba firmado o no pudo firmarse después. Desde la fecha cierta, el instrumento privado es \textbf{oponible a terceros}. El CCyCN abandona la enumeración taxativa del Código anterior: la fecha cierta puede acreditarse por cualquier medio, que el juez apreciará rigurosamente. \\[6pt]

\textbf{¿Cómo valora el juez un instrumento particular no firmado? (Art. 319)} &
Debe ponderar: congruencia entre lo narrado y lo ocurrido; precisión y claridad técnica del texto; usos y prácticas del tráfico; relaciones precedentes entre las partes; y confiabilidad de los soportes y procedimientos técnicos empleados. \\[6pt]

\textbf{¿Puede utilizarse la correspondencia como prueba en juicio? (Art. 318)} &
Sí: el destinatario puede presentarla como prueba contra el remitente. Sin embargo, si la correspondencia es \textbf{confidencial} —por manifestación expresa o por la naturaleza de la comunicación—, no puede utilizarse sin consentimiento del remitente. Esta protección tiene fundamento constitucional (inviolabilidad de la correspondencia). \\

\midrule

\multicolumn{2}{p{\textwidth}}{%
\textbf{Resumen:} La forma del acto jurídico es el modo en que la voluntad de las partes se exterioriza y produce efectos jurídicos. El CCyCN parte del principio de libertad de formas (Art. 284), pero impone formas determinadas para los actos más trascendentes, con distintas consecuencias según se trate de solemnidad absoluta (nulidad del acto) o relativa (conversión: obligación de otorgar la forma prescripta, Art. 285). En materia de instrumentos, el Código distingue tres grandes categorías: los \textbf{instrumentos públicos}, que gozan de autenticidad propia y hacen plena fe, con dos regímenes de impugnación según se trate de hechos percibidos por el oficial (redargución de falsedad) o de declaraciones de las partes (prueba en contrario); la \textbf{escritura pública}, instrumento matriz del protocolo notarial que contiene actos jurídicos y requiere una estructura formal precisa; y los \textbf{instrumentos particulares}, que incluyen los privados (requisito esencial: firma; fuerza probatoria plena tras el reconocimiento) y los particulares no firmados (fuerza probatoria sujeta a la apreciación judicial según las pautas del Art. 319). La \textbf{firma digital} (Ley 25.506) garantiza autoría e integridad en el mundo digital, con mayor valor probatorio que la firma electrónica. La \textbf{fecha cierta} (Art. 317) determina desde cuándo el instrumento privado puede oponerse a terceros, y puede acreditarse por cualquier medio apreciado rigurosamente por el juez. Estas reglas configuran un sistema orientado a dar seguridad jurídica al tráfico de bienes y servicios, equilibrando la certeza del formalismo con la agilidad que demanda la vida contemporánea.
}\\

\bottomrule
\end{tabularx}
\end{table}

% ============================================================
\end{document}
% ============================================================
